\section{August 31, 2026}

We continue the analysis of one-step methods. The main goals are to pass
from local consistency to global convergence and to introduce
Runge--Kutta methods. Consider the IVP
\[
  x'(t)=f(t,x(t)),
  \qquad
  x(t_0)=x_0\in\bb{R}^n,
\]
on a grid
\[
  t_0<t_1<\cdots<t_N=T,
  \qquad
  \tau_j=t_{j+1}-t_j,
  \qquad
  h=\max_{0\leq j<N}\tau_j.
\]

\subsection{From local consistency to global convergence}

Recall that \(\Phi^{t_{j+1},t_j}\) denotes the exact flow over one
step, while \(\Psi^{t_{j+1},t_j}\) denotes the numerical one-step map.
For a starting value \(x\), define the local consistency error by
\begin{equation}\label{eq:local-error-aug-31}
  \xi_j(x)
  =\Phi^{t_{j+1},t_j}(x)
   -\Psi^{t_{j+1},t_j}(x).
\end{equation}

\begin{definition}[Consistency order]
  A one-step method has \emph{consistency order} \(p>0\) if, uniformly
  for the relevant starting values,
  \[
    \lVert\xi_j(x)\rVert\leq C\tau_j^{p+1}
  \]
  as \(h\to0\). The order \(p\) need not be an integer when the data
  have only H\"older regularity.
\end{definition}

Forward and backward Euler have consistency order one. The exact and
numerical values at the final time can be written as compositions:
\begin{align*}
  x(T)
  &=\Phi^{t_N,t_{N-1}}\circ\cdots\circ
    \Phi^{t_2,t_1}\circ\Phi^{t_1,t_0}(x_0),\\
  x_N
  &=\Psi^{t_N,t_{N-1}}\circ\cdots\circ
    \Psi^{t_2,t_1}\circ\Psi^{t_1,t_0}(x_0).
\end{align*}
Thus the local errors are propagated through all the subsequent steps.

Let
\[
  e_j=x_j-x(t_j)
\]
be the global error. For brevity, set
\(\Psi_j=\Psi^{t_{j+1},t_j}\) and
\(\Phi_j=\Phi^{t_{j+1},t_j}\). Then
\begin{align}
  e_{j+1}
  &=\Psi_j(x_j)-\Phi_j(x(t_j))\notag\\
  &=\Psi_j(x_j)-\Psi_j(x(t_j))-\xi_j(x(t_j)).
  \label{eq:error-recursion-aug-31}
\end{align}
Suppose the one-step map satisfies the stability estimate
\begin{equation}\label{eq:step-stability-aug-31}
  \lVert\Psi_j(y)-\Psi_j(z)\rVert
  \leq e^{L\tau_j}\lVert y-z\rVert.
\end{equation}
Combining the consistency and stability estimates in
\eqref{eq:error-recursion-aug-31} gives
\begin{equation}\label{eq:error-bound-one-step-aug-31}
  \lVert e_{j+1}\rVert
  \leq e^{L\tau_j}\lVert e_j\rVert+C\tau_j^{p+1}.
\end{equation}
Since \(e_0=0\), iteration yields
\begin{align*}
  \lVert e_N\rVert
  &\leq C\sum_{j=0}^{N-1}
    \tau_j^{p+1}e^{L(T-t_{j+1})}\\
  &\leq Ch^p\sum_{j=0}^{N-1}
    \tau_j e^{L(T-t_{j+1})}.
\end{align*}
For \(L>0\), the sum is bounded by a geometric integral, and hence
\begin{equation}\label{eq:global-error-bound-aug-31}
  \lVert e_N\rVert
  \leq \frac{C}{L}\bigl(e^{L(T-t_0)}-1\bigr)h^p.
\end{equation}
When \(L=0\), the corresponding bound is
\[
  \lVert e_N\rVert\leq C(T-t_0)h^p.
\]
In either case, \(h\to0\) implies \(e_N\to0\). Thus stability converts
a local error of order \(p+1\) into a global error of order \(p\).

\begin{remark}[Contractive problems]
  In the favorable case
  \[
    \lVert\Psi_j(y)-\Psi_j(z)\rVert\leq\lVert y-z\rVert,
  \]
  no exponential amplification occurs. Direct summation gives
  \[
    \lVert e_N\rVert
    \leq C\sum_{j=0}^{N-1}\tau_j^{p+1}
    \leq C(T-t_0)h^p.
  \]
  Such contractivity is especially relevant for dissipative ODEs and
  numerical methods that preserve their dissipative structure.
\end{remark}

\subsection{Runge--Kutta methods}

Consider an explicit method with \(s\) stages and coefficients
\(a_{ij}\), \(b_i\), and \(c_i\). Starting from \(x\) at time \(t\)
with step size \(\tau\), its stages are computed successively from
\begin{equation}\label{eq:rk-stages-aug-31}
  k_i=f\left(t+c_i\tau,
    x+\tau\sum_{j=1}^{i-1}a_{ij}k_j\right),
  \qquad i=1,\ldots,s,
\end{equation}
and its one-step map is
\begin{equation}\label{eq:rk-map-aug-31}
  \Psi^{t+\tau,t}(x)
  =x+\tau\sum_{i=1}^s b_i k_i.
\end{equation}

\begin{example}[Forward Euler]
  Take \(s=1\), \(c_1=0\), \(a_{11}=0\), and \(b_1=1\). Then
  \[
    k_1=f(t,x),
    \qquad
    \Psi^{t+\tau,t}(x)=x+\tau k_1,
  \]
  which is forward Euler, a first-order method.
\end{example}

\begin{example}[Explicit trapezoidal method]
  Take \(s=2\) and
  \[
    A=\begin{bmatrix}0&0\\1&0\end{bmatrix},
    \qquad
    b=\begin{bmatrix}\tfrac12\\[2pt]\tfrac12\end{bmatrix},
    \qquad
    c=\begin{bmatrix}0\\1\end{bmatrix}.
  \]
  The stages and update are
  \begin{align*}
    k_1&=f(t,x),\\
    k_2&=f(t+\tau,x+\tau k_1),\\
    \Psi^{t+\tau,t}(x)&=x+\frac{\tau}{2}(k_1+k_2).
  \end{align*}
  This is the predictor--corrector method from the previous lecture,
  also called Heun's method. It has order two.
\end{example}
